% ======================================================================
\section{Error analysis}\label{sec:analysis}

This section looks at the remaining errors of the final system in detail, using its logged decision trail, per-item
error lists and an audit of the on-screen check (sources in Appendix~\ref{app:sources}).

\subsection{Ceiling analysis}\label{sec:ceiling}
Table~\ref{tab:ceiling} gives the numbers behind Figure~\ref{fig:ceiling}: each row makes one stage perfect on its own.
The analysis needs saved pictures, so it was run on a slightly older version of the system (25 hits and 27 wrong on the
development set, 55 needed sounds); the proportions, not the exact numbers, are the point.

\begin{table}[h]
\centering\small
\caption{One stage made perfect at a time (development set, slightly older version).}
\label{tab:ceiling}
\begin{tabular}{lrrr}
\toprule
Stage made perfect & Hits & Wrong & Cost (from 2.451) \\
\midrule
On-screen check & 30 & 18 & 1.915 \\
Listeners and their filters & 34 & 24 & 1.859 \\
Vetoes and filters & 30 & 27 & 2.169 \\
Timing & 26 & 20 & 2.197 \\
\midrule
All stages perfect & 43 & 3 & 0.761 \\
\bottomrule
\end{tabular}
\end{table}

The on-screen check and the listeners are the two largest losses. Even with every stage perfect, 12 of 55 needed
sounds stay missed: five are heard by no model at all (a hammer, an explosion, a clang and two golf-club strikes) and
seven are too faint or too badly timed.

\subsection{The on-screen check in detail}\label{sec:presence}
The audit of the gate's decisions found two kinds of error. The more common one (23 of 36 gate errors) is
\emph{presence taken for production}: an object in the frame is taken as the object making the sound. This follows from
how the vote is built. Two of the three questions effectively test presence -- naming an object, and asking whether that
object \emph{could} make the sound -- and only the a/b question asks whether the sound is visibly happening, so the
majority tests presence too. The second error is the opposite: the name question answers ``nothing'' on 42\,\% of the
stretches the labels call visible, when the source is small, hidden, or its action is fine.

\subsection{Label strictness}
In two blind checks, about 1 in 6 pictures counted as wrong was a real, wanted, off-screen sound that the labels do not
list, and several others were real sounds with the source on screen. Absolute costs are therefore somewhat pessimistic;
the comparison between systems is not affected, because the same labels score all of them.

\subsection{Worked example}
Clip \fp{bell_miami} (development set) has one needed sound: a bell from 0.2 to 14.5\,s, labelled not visible. The
detectors find it. The gate cuts the 14-s sound into three stretches. In each, the name question answers ``church
bell'' (it shares the word \emph{bell}, so this vote says visible), the a/b question says not visibly happening, and the
description vote says visible. Two of three votes in every stretch: the sound is silenced and the needed bell is missed.
The frames show the church; they do not show the bell ringing. The one question that asked about the action was right,
and it was outvoted by two questions that measure presence.
